\documentclass[12pt]{article}
\usepackage[utf8]{inputenc}
\usepackage[T1]{fontenc}
\usepackage{hyperref}
\title{Minidien Language Resources for K-12 Curriculum Development}
\author{Compiled by Claude Sonnet 5.5 (claude-sonnet-5-5) for Andr\'as Kornai}
\date{}
\begin{document}
\maketitle

\section{Introduction}
Minidien (also Wiaki, Miniden, Menandon; ISO 639-3 wii) is a Torricelli language of Sandaun Province, Papua New Guinea. Sources disagree on speaker numbers: about 100 (2004) on Wikipedia and about 200 on Joshua Project, with an ethnic population of about 1{,}190 (2000 census). Use is concentrated among older people. The web search found very little beyond encyclopedic overviews.

\section{Primary Educational Resources}
No primer, textbook, or school programme for Minidien turned up in search. None should be assumed to exist.

\section{Linguistic and Reference Materials}
Overview of the language and its place in the Maimai branch of Torricelli: \url{https://en.wikipedia.org/wiki/Minidien_language}. Family context: \url{https://en.wikipedia.org/wiki/Torricelli_languages}. No dictionary or grammar was found.

\section{Religious and Community Resources}
Joshua Project has a people-group profile (as Wiaki): \url{https://joshuaproject.net/people_groups/15396}. Its Bible translation status is unspecified. No scripture was found.

\section{Digital and Audio Resources}
No audio, video, or corpus was found in the search results.

\section{Curriculum Development Context}
Documentation is sparse, so a curriculum would have to start from new documentation work with the community and older speakers. Community organisations and any existing fieldwork archives need to be identified; the search did not find them.
\end{document}